% Shared body for the inverter wiring figures. The wrappers diagram_gfl.tex
% and diagram_gfm.tex set \gfmfalse or \gfmtrue before inputting this file;
% the switch selects the outer controller, its setpoint inputs, and the extra
% voltage and frequency wiring the voltage loop needs.

% TUNABLES: tier heights, the PLL frame boundary, and the setpoint column.
\def\framerow{-2.8}
\def\framecolumn{6.8}
\def\controlrow{-4.4}
\def\setpointcolumn{18.6}

\begin{document}
\begin{tikzpicture}[emt diagram,
    block/.append style={minimum height=1.25cm, align=center}]

% ================= DC supply and bridge =================
\node[signal, label={[sig]above:$v_{\mathrm{dc}}$}] (vdc) at (4.8, 0) {};
\node[block, minimum width=2.6cm] (bridge) at (8.8, 0) {Converter};
\node[signal] (dcbranch) at (7.2, 0) {};
\draw[->, line] (vdc) -- (bridge.west);

% ================= LCL filter and terminal =================
\node[block] (filter) at (12.4, 0) {Filter\\{\small LCL}};
\node[block] (tbus) at (20.0, 0) {Bus\\{\small terminal}};
\draw[->, line] (bridge.east)
    -- node[sig, above] {$\mathbf{e}$} (filter.west);
\draw[->, line] (filter.east) -- (tbus.west);
% The grid-current tap sits over the outer controller input it feeds.
\node[signal, label={[sig]above:$\mathbf{i}_g$}] (igabc)
    at (\ifgfm 16.15\else 15.9\fi, 0) {};

% ================= Reference frame and control =================
\node[block] (icc) at (12.4, \controlrow) {ICC};
\node[block] (outer) at (15.9, \controlrow) {\ifgfm OVC\else OPC\fi};

% ================= Phase-domain measurements =================
% The filter publishes its inverter current and capacitor voltage from below;
% reference-frame transformations are implicit at the dq boundary.
\draw[->, line] ([xshift=-2.5mm]filter.south) -- ([xshift=-2.5mm]icc.north);
\node[sig, anchor=east] at (12.0, -1.2) {$\mathbf{i}$};
\draw[->, line] ([xshift=2.5mm]filter.south) -- ([xshift=2.5mm]icc.north);
\node[sig, anchor=west] at (12.8, -1.2) {$\mathbf{v}_{\mathrm{o}}$};
\draw[->, line] (igabc) -- ([xshift=\ifgfm 2.5mm\else 0mm\fi]outer.north);

\ifgfm
% ================= Outer voltage-control wiring =================
% The voltage loop also reads the capacitor voltage.
\node[signal] (vouter) at (12.65, -2.2) {};
\draw[->, line] (vouter) -- (15.65, -2.2) -- ([xshift=-2.5mm]outer.north);
\fi

% ================= Control chain =================
\draw[->, line] ([yshift=-2.5mm]outer.west)
    -- node[sig, below] {$\mathbf{i}^{\mathrm{cmd}}$} ([yshift=-2.5mm]icc.east);
\draw[->, line] ([yshift=2.5mm]icc.east)
    -- node[sig, above] {$\mathbf{i}^{\mathrm{lim}}$} ([yshift=2.5mm]outer.west);

% ================= PWM =================
% The modulator sits under the bridge inside the dq frame: it limits the
% voltage command, returns the limited command, and sends the switching
% functions straight up to the bridge.
\node[block] (pwm) at (8.8, \controlrow) {PWM};
\draw[->, line] (pwm.north) -- node[sig, right] {$\mathbf{s}$} (bridge.south);
\draw[->, line] ([yshift=2.5mm]icc.west)
    -- node[sig, above] {$\mathbf{u}$} ([yshift=2.5mm]pwm.east);
\draw[->, line] ([yshift=-2.5mm]pwm.east)
    -- node[sig, below] {$\mathbf{u}^{\mathrm{lim}}$} ([yshift=-2.5mm]icc.west);
\draw[->, line] (dcbranch) |- (pwm.west);

% ================= Terminal-voltage PLL =================
\node[block] (pll) at (20, \framerow) {PLL};
\draw[->, line] (tbus.south) -- node[sig, right] {$\mathbf{v}$} (pll.north);

% ================= Setpoints =================
% Setpoints enter the outer controller from outside the enclosure.
\ifgfm
\node[signal, label={[sig]above:$\mathbf{v}^{\mathrm{ref}}$}] (vref) at (\setpointcolumn, \controlrow) {};
\draw[->, line] (vref) -- (outer.east);
\else
\node[signal, label={[sig]above:$P^{\mathrm{ref}}$}] (pref) at (\setpointcolumn, -4.15) {};
\draw[->, line] (pref) -- ([yshift=2.5mm]outer.east);
\node[signal, label={[sig]below:$Q^{\mathrm{ref}}$}] (qref) at (\setpointcolumn, -4.65) {};
\draw[->, line] (qref) -- ([yshift=-2.5mm]outer.east);
\fi

% ================= Configuration and frame enclosures =================
\begin{scope}[on background layer]
    \node[operator enclosure, densely dotted,
        fit=(bridge)(filter)
            (icc)(outer)(pwm)] (enclosure) {};
    % The PLL defines the dq frame; the region is labelled with it.
    \coordinate (frame) at (\framecolumn, \framerow);
    \draw[enclosure, densely dotted]
        (frame |- enclosure.south) -- (frame) -- (frame -| enclosure.east);
\end{scope}
\draw[->, line] (pll.west)
    -- node[sig, above] {$\theta,\omega$} (frame -| enclosure.east);
\node[sig, anchor=north west] at ($(frame)+(0.7,-0.1)$) {$dq$};

\end{tikzpicture}
\end{document}
